\section{Hệ thống đề xuất}
\subsection{Mô tả hệ thống}
Hệ thống được thiết kế nhằm hỗ trợ phân tích và đánh giá tình trạng giao thông (TTGT) trong môi trường đô thị thông qua việc xử lý, tổ chức và phân tích dữ liệu giao thông được cung cấp. Hệ thống cho phép mô hình hóa mạng lưới giao thông, nhận diện các mối quan hệ giữa các đoạn đường và khai thác dữ liệu theo cả hai khía cạnh không gian và thời gian để hỗ trợ đánh giá và hiểu rõ hơn về tình trạng lưu thông.

\subsubsection{Mục đích của hệ thống}
Mục tiêu chính của hệ thống bao gồm:
\begin{itemize}
    \item Quản lý dữ liệu giao thông từ nhiều nguồn (dữ liệu bản đồ số, dữ liệu cảm biến, dữ liệu GPS, dữ liệu từ các nền tảng giao thông của bên thứ ba, v.v.).
    \item Tiền xử lý, chuyển đổi và chuẩn hóa dữ liệu để phục vụ phân tích.
    \item Mô hình hóa mối liên hệ giữa các đoạn đường trong mạng lưới giao thông nhằm hiểu rõ cấu trúc và dòng chảy giao thông.
    \item Phân tích dữ liệu để hỗ trợ nhận diện tình trạng ùn tắc, đánh giá xu hướng và hỗ trợ ra quyết định.
    \item Trực quan hóa kết quả phân tích bằng biểu đồ, bản đồ hoặc báo cáo nhằm giúp người dùng theo dõi và hiểu kết quả.
    \item Cung cấp nền tảng linh hoạt để thử nghiệm, so sánh hoặc tích hợp các phương pháp phân tích trong tương lai.
\end{itemize}

\subsubsection{Nhóm người dùng mục tiêu}
Hệ thống hướng đến ba nhóm người dùng chính:
\begin{itemize}
    \item \textbf{Nhà nghiên cứu và sinh viên:} sử dụng hệ thống để phân tích dữ liệu giao thông, thử nghiệm các phương pháp mô hình hóa khác nhau và đánh giá các yếu tố ảnh hưởng đến TTGT.
    \item \textbf{Cơ quan quản lý giao thông hoặc nhà quy hoạch đô thị:} khai thác hệ thống để theo dõi tình trạng giao thông, phát hiện điểm nóng ùn tắc và đánh giá tác động của các giải pháp tổ chức giao thông.
    \item \textbf{Kỹ sư dữ liệu và nhà phát triển hệ thống:} sử dụng hệ thống như nền tảng để xử lý dữ liệu giao thông, mở rộng mô hình, tích hợp thêm tính năng hoặc triển khai vào môi trường thực tế.
\end{itemize}

\subsubsection{Tổng quan}
Tổng thể, hệ thống được thiết kế theo hướng mở, linh hoạt. Điều này cho phép hệ thống dễ dàng mở rộng, thay đổi hoặc kết hợp các phương pháp phân tích khi cần thiết, đồng thời đảm bảo phù hợp cho cả mục tiêu nghiên cứu lẫn ứng dụng thực tiễn trong phân tích tình trạng giao thông đô thị.
\subsection{Yêu cầu}
\subsubsection{Yêu cầu chức năng}

\paragraph{Tiền xử lý dữ liệu}
\begin{itemize}
    \item \textbf{FR1.1}: Loại bỏ dữ liệu thiếu, sai định dạng hoặc bất thường.
    \item \textbf{FR1.2}: Chuẩn hoá các thuộc tính như tốc độ, thời gian di chuyển, timestamp.
    \item \textbf{FR1.3}: Ghép nối dữ liệu với mạng lưới đường (map-matching) khi cần.
\end{itemize}

\paragraph{Phân tích mạng lưới giao thông}
\begin{itemize}
    \item \textbf{FR2.1}: Xây dựng mô hình đồ thị gồm các nút (node) và cạnh (segment).
    \item \textbf{FR2.2}: Tính toán các thông số của mạng đường: độ dài, phân cấp đường (FRC), các chỉ số trung tâm.
    \item \textbf{FR2.3}: Xác định “điểm nóng” (hotspot) tắc nghẽn theo thời gian.
\end{itemize}

\paragraph{Phân tích mối liên hệ giữa các đoạn đường}
\begin{itemize}
    \item \textbf{FR3.1}: Tính toán mức độ tương quan giữa các đoạn đường theo thời gian.
    \item \textbf{FR3.2}: Đo độ tương đồng dạng thời gian bằng kỹ thuật Dynamic Time Warping (DTW).
    \item \textbf{FR3.3}: Phát hiện hiện tượng lan truyền tắc nghẽn trong mạng lưới.
\end{itemize}

\paragraph{Dự báo lan truyền ùn tắc giao thông}
\begin{itemize}
    \item \textbf{FR4.1}: Xây dựng tập dữ liệu theo chuỗi thời gian kết hợp với cấu trúc mạng lưới giao thông nhằm mô hình hóa sự lan truyền của ùn tắc giữa các nút và đoạn đường.
    \item \textbf{FR4.2}: Áp dụng và huấn luyện các mô hình dự báo lan truyền ùn tắc giao thông trên mạng lưới giao thông.
    \item \textbf{FR4.3}: Đánh giá hiệu quả và độ chính xác của mô hình dự báo lan truyền thông qua các thước đo như MAE, RMSE, đồng thời phân tích khả năng phản ánh xu hướng lan truyền theo không gian và thời gian.
\end{itemize}

\paragraph{Trực quan hoá dữ liệu}
\begin{itemize}
    \item \textbf{FR5.1}: Hiển thị bản đồ mạng giao thông bằng dữ liệu OSM hoặc TomTom.
    \item \textbf{FR5.2}: Trực quan tốc độ trung bình theo thời gian của từng segment.
    \item \textbf{FR5.3}: Hiển thị biểu đồ chuỗi thời gian.
    \item \textbf{FR5.4}: Hiển thị ma trận tương quan giữa các đoạn đường.
    \item \textbf{FR5.5}: Cho phép thao tác xem dữ liệu trên dashboard.
\end{itemize}

\paragraph{Xuất báo cáo}
\begin{itemize}
    \item \textbf{FR6.1}: Xuất kết quả phân tích ra các định dạng CSV, JSON hoặc PDF.
    \item \textbf{FR6.2}: Cho phép tải xuống kết quả.
\end{itemize}
\subsubsection{Yêu cầu phi chức năng}
\paragraph{Khả năng mở rộng}
\begin{itemize}
    \item Hệ thống phải có khả năng xử lý dữ liệu lớn lên đến hàng triệu bản ghi, đồng thời có thể mở rộng mà không suy giảm hiệu năng.
\end{itemize}

\paragraph{Hiệu năng}
\begin{itemize}
    \item Thời gian xử lý dữ liệu phải đảm bảo hiệu quả.
    \item Thời gian tải và hiển thị biểu đồ phải dưới 5 giây.
\end{itemize}

\paragraph{Độ tin cậy}
\begin{itemize}
    \item Hệ thống phải hoạt động ổn định và cho kết quả nhất quán khi sử dụng cùng một bộ dữ liệu.
\end{itemize}

\paragraph{Tính bảo trì}
\begin{itemize}
    \item Mã nguồn phải được tổ chức thành các module độc lập gồm: tiền xử lý, phân tích, mô hình hoá và trực quan hoá.
\end{itemize}

\paragraph{Tính tương thích}
\begin{itemize}
    \item Hệ thống phải hoạt động tốt với Python 3.10+, các thư viện phân tích (Pandas, GeoPandas, NetworkX), TomTom MOVE và dữ liệu OSM.
\end{itemize}
\paragraph{Tính dễ sử dụng}
\begin{itemize}
    \item Giao diện và quy trình thao tác phải đơn giản, giúp người dùng dễ dàng tải dữ liệu, xem biểu đồ và tạo báo cáo.
\end{itemize}
\paragraph{An toàn dữ liệu}
\begin{itemize}
    \item Không thu thập hoặc xử lý dữ liệu định danh cá nhân. Dữ liệu giao thông chỉ được xử lý ở mức tổng hợp.
\end{itemize}
\subsubsection{Yêu cầu dữ liệu}

\subsubsubsection{Loại dữ liệu cần thu thập}
Hệ thống cần dữ liệu mô tả tình trạng giao thông đô thị có thể từ một hoặc nhiều nguồn khác nhau (API giao thông, cảm biến, dữ liệu mở, dữ liệu bản đồ số, v.v.). Các nhóm dữ liệu chính bao gồm:
\begin{itemize}
    \item \textbf{Dữ liệu cấu trúc mạng lưới giao thông:}
    \begin{itemize}
        \item Thông tin các đoạn đường (segment).
        \item Mã định danh đoạn đường.
        \item Tên đường, khu vực, vùng hoặc zone.
        \item Thông tin bản đồ nền và phiên bản bản đồ.
    \end{itemize}

    \item \textbf{Dữ liệu thời gian:}
    \begin{itemize}
        \item Thời điểm thu thập dữ liệu.
        \item Khoảng ngày quan sát.
        \item Khoảng thời gian trong ngày.
    \end{itemize}

    \item \textbf{Dữ liệu tình trạng giao thông:}
    \begin{itemize}
        \item Tốc độ trung bình, tốc độ tự do.
        \item Mức độ chậm trễ, chỉ số ùn tắc.
        \item Điều kiện lưu thông trên từng đoạn đường.
        \item Nguồn thu thập (GPS, cảm biến, v.v.) nếu có.
    \end{itemize}

    \item \textbf{Thông tin cấu hình và siêu dữ liệu:}
    \begin{itemize}
        \item Tên phiên thu thập hoặc mã job.
        \item Đơn vị đo vận tốc.
        \item Các lựa chọn cấu hình do người dùng đặt.
        \item Các tham số được sử dụng để tạo bộ dữ liệu.
    \end{itemize}
\end{itemize}

\subsubsubsection{Yêu cầu về định dạng và lưu trữ dữ liệu}
\begin{itemize}
    \item Dữ liệu đầu vào có thể ở dạng JSON, CSV hoặc các định dạng bán cấu trúc khác, tùy theo nguồn thu thập.
    \item Hệ thống cần hỗ trợ chuyển đổi sang dạng bảng (tabular) để phân tích.
    \item Dữ liệu được lưu trữ ở các định dạng thuận tiện cho xử lý như CSV, Parquet...
    \item Hỗ trợ nhập dữ liệu từ nhiều nguồn: file tải lên, API hoặc kho lưu trữ trên nền tảng đám mây.
\end{itemize}

\subsubsubsection{Yêu cầu về chất lượng dữ liệu}
\begin{itemize}
    \item \textbf{Đầy đủ:} có đủ thông tin tối thiểu để phân tích tình trạng giao thông.
    \item \textbf{Nhất quán:} đồng nhất về đơn vị đo, định dạng thời gian và mã định danh.
    \item \textbf{Đúng cấu trúc:} các trường phải đúng kiểu dữ liệu.
    \item \textbf{Hạn chế dữ liệu lỗi hoặc thiếu:} hạn chế phần tử bị trống hoặc sai lệch.
    \item \textbf{Có thể truy vết:} biết rõ nguồn và thời điểm thu thập.
\end{itemize}

\subsubsubsection{Yêu cầu về cập nhật dữ liệu}
\begin{itemize}
    \item Hệ thống cần hỗ trợ cập nhật dữ liệu từ nhiều nguồn khác nhau trong tương lai.
    \item Hệ thống cho phép cập nhật định kỳ khi có dữ liệu mới.
\end{itemize}